
Foundation models transformed not only AI capabilities but also how researchers evaluate progress. This study presents quantitative evidence that foundation model emergence coincided with a phase transition in ML benchmark concentration dynamics. Applying PELT change-point detection to 84 months of Papers With Code data (2018--2024), two structural breaks were identified---April 2019 and March 2021---with a BIC improvement of 17.13 over a monotonic-trend null hypothesis. The underlying mechanism is attention reallocation: emergent-capability benchmarks attracted +19 percentage points more researcher attention post-2021, while traditional benchmarks persisted with reduced dominance. Contrary to initial expectations, modality-specific dynamics (CV vs. NLP) were never unified pre-2020 (Pearson r = $-0.13)$, indicating that foundation models restructured rather than fragmented the benchmark ecosystem. Five of six falsifiable hypotheses were validated (83.3\%), with both MUST\_WORK gates passing. This hypothesis-driven methodology provides a template for rigorous meta-science of AI research practices.
